\chapter{Despliegue, respaldo y recuperación}

\section{Inventario operativo}

Antes de transferir la operación, registre para cada servicio: propietario, segundo administrador,
facturación, dominio, método de recuperación y ubicación segura de secretos.

\begin{tabularx}{\textwidth}{@{}L{0.22\textwidth}Y Y@{}}
\toprule
\textbf{Servicio} & \textbf{Responsabilidad} & \textbf{Verificación mínima}\\
\midrule
GitHub & código, PR, CI y releases & ownership, protección de `main`, tag final\\
Vercel & frontend y rewrite `/api` & build, variables, dominio, ruta profunda\\
Render & API/ASGI y logs & health, comando, variables, rollback\\
Supabase & PostgreSQL y Storage & backups, acceso, bucket y límites\\
Redis & channel layer & conectividad, aislamiento y persistencia requerida\\
Correo & Brevo o SMTP & remitente validado y entrega controlada\\
\bottomrule
\end{tabularx}

\section{Variables}

Las plantillas canónicas son \path{backend/.env.example}, \path{frontend/.env.example} y
\path{docker-compose.prod.env.example}. Nunca copie un archivo `.env` real al repositorio.

\begin{itemize}
  \item Backend base: \path{DJANGO_SECRET_KEY}, \path{DJANGO_DEBUG},
  \path{DATABASE_URL} y \path{ALLOWED_HOSTS}.
  \item Orígenes y sesión: \path{CORS_ALLOWED_ORIGINS},
  \path{CSRF_TRUSTED_ORIGINS}, \path{FRONTEND_URL},
  \path{AUTH_COOKIE_SAMESITE} y \path{AUTH_COOKIE_SECURE}.
  \item Frontend: \path{VITE_API_BASE_URL}, \path{VITE_WS_URL},
  \path{VITE_ENV} y la opción de transformación de imágenes si se utiliza.
\end{itemize}

\section{Release controlada}

\begin{enumerate}
  \item Confirme PR aprobado, CI verde, SHA y notas de cambio.
  \item Revise `showmigrations` y `migrate --plan`.
  \item Cree un respaldo válido antes de una migración con riesgo de datos.
  \item Ejecute migraciones una sola vez como paso de release.
  \item Despliegue backend y frontend desde referencias inmutables.
  \item Ejecute health y el smoke test por roles.
  \item Cree tag/release y registre el resultado o active rollback.
\end{enumerate}

\begin{lstlisting}[language=bash]
cd backend
python manage.py showmigrations
python manage.py migrate --plan
python manage.py migrate
python manage.py check --deploy
\end{lstlisting}

\RiskWarning{No ejecute `flush`, no borre migraciones aplicadas y no use `seed\_demo` contra
producción. El arranque normal del contenedor no debe ejecutar migraciones.}

\section{Smoke test}

Compruebe, en orden: frontend y ruta profunda; `/health/`; inicio de sesión controlado; creación de
ticket; asignación; comentario y cambio de estado; actualización del Cliente; reporte PDF/Excel;
WebSocket con respuesta 101; y correo desde el remitente esperado.

Registre fecha, SHA/tag, entorno, responsable y resultado. Use datos ficticios y elimine evidencia
con datos personales antes de compartirla.

\section{Backup}

Priorice los backups administrados y PITR de Supabase cuando el plan los incluya. Para una copia
manual:

\begin{enumerate}
  \item ejecute `pg\_dump` desde un host autorizado;
  \item valide que `pg\_restore` pueda leer el dump;
  \item calcule SHA-256 y registre fecha, tamaño, retención y responsable;
  \item almacene el archivo cifrado y separado de las credenciales;
  \item inventaríe los objetos de Storage por separado.
\end{enumerate}

El script `scripts/backup\_db.sh` automatiza dump, validación estructural y checksum. No demuestra
por sí mismo que el sistema pueda restaurarse.

\section{Restauración y rollback}

Restaure primero en una base aislada. Ejecute checks y migraciones, compare conteos y relaciones,
pruebe autenticación y un ciclo completo, confirme Storage y registre los RTO y RPO observados. Ante
una release fallida, detenga el rollout, preserve logs, revierta a la última referencia conocida,
aplique solo el plan de datos aprobado, repita health check y smoke tests y rote secretos si el
incidente fue de seguridad.
